\documentclass[a4paper,10pt]{article}
\usepackage{amsmath}
\usepackage[utf8]{inputenc}
\usepackage{geometry}
\usepackage{palatino}
\geometry{margin=1cm}

\begin{document}
\centerline{\textbf{Evaluación: Termodinámica} \quad Alumno: \underline{\hspace{4cm}} \quad Escuela: 4-117 Ejército de los Andes}

1) Juan compró \emph{Termotanque Eléctrico Señorial Safirus 65} por Mercadolibre, al llegar por \emph{Correo Andreani} a su domicilio e instalarlo el plomero Alberto, pudo ver que a pesar de ser 65 posee un tanque de $V=50L$ ya está escrito en la etiqueta de especificaciones técnicas, además tiene una resistencia para calentar el agua consume P=1,5kW, el agua fría $Cm_{H_2O}=1 kcal/{kg^\circ C}$  ingresa a $t_1=20 ^\circ F$ de la red, suponiendo que la aislación térmica es perfecta, si se desea que el agua llegue a $t_2=45°C$ para bañarse con agua caliente.  ¿Cuánto calor en kcal deberá transferir la resistencia al agua? 
 
 \vspace{0.5cm}

2) Carlos compró una \emph{Heladera No Frost Whirlpool WRE57K3 438L} en MercadoLibre. Supongamos que la temperatura ambiente es de $30^\circ C$ y el interior de la heladera debe mantenerse a $4^\circ F$.  
Las paredes tienen un espesor de $5cm$ y están hechas de un material con una conductividad térmica de $0.025 \frac{J}{m \cdot K \cdot s}$.  
Si la superficie total de la heladera es $2.5 m^2$, ¿cuánto calor se transfiere por hora a través de las paredes de la heladera?
  

\vspace{0.5cm}

3)
\begin{itemize}
    \item $127^\circ K \xrightarrow{\text{convertir}}  .^\circ F $

\item $362^\circ F \xrightarrow{\text{convertir}}  .^\circ K $

\item $0,18^\circ Chu \xrightarrow{\text{convertir}} BTU $

\item $10,14^\circ kcal \xrightarrow{\text{convertir}}  BTU $
 
\end{itemize}
    



\rule{\textwidth}{1pt}    
\centerline{\textbf{Evaluación: Termodinámica} \quad Alumno: \underline{\hspace{4cm}} \quad Escuela: 4-117 Ejército de los Andes}

1) Ana compró un \emph{Hervidor Eléctrico Philips Daily Collection 1.7L 2200W} en Amazon. Quiere hervir $1.5L$ de agua que está a $t_1 = 18^\circ C$ hasta alcanzar los $300^\circ K$. Dado que el calor específico del agua es $1 \text{ kcal/kg}^\circ C$ y suponiendo que no hay pérdidas térmicas:   ¿Cuánto calor en kcal debe transferir la resistencia al agua?



\vspace{0.5cm}
  

2) Martín tiene una \emph{Pava Eléctrica Atma 1.7L 2000W}. La base de la pava es de \emph{acero inoxidable} con un espesor de $3mm$ y una conductividad térmica de $16 \frac{W}{m \cdot ^\circ C}$.  
Si el área de contacto con el agua es $0.02m^2$, y la diferencia de temperatura entre el agua y la resistencia es de $380^\circ K$, ¿Cuánto calor se transfiere por segundo a través de la base de la pava?

3)
\begin{itemize}
  
\item $36^\circ F \xrightarrow{\text{convertir}}  .^\circ K $

\item $0,18^\circ BTU \xrightarrow{\text{convertir}} kcal $

\item $1144^\circ BTU \xrightarrow{\text{convertir}}  kcal $

   \item $128^\circ K \xrightarrow{\text{convertir}}  .^\circ F $

\end{itemize}

\rule{\textwidth}{1pt} 
\centerline{\textbf{Evaluación: Termodinámica} \quad Alumno: \underline{\hspace{4cm}} \quad Escuela: 4-117 Ejército de los Andes}

1) Marcos compró un \emph{Termotanque Eléctrico Rheem 85L 2000W} en MercadoLibre.  
El agua de la red entra a $t_1 = 15^\circ C$ y él quiere calentar toda el agua a $t_2 = 50^\circ C$ para una ducha.  
Dado que el calor específico del agua es $1 \text{ kcal/kg}^\circ C$ y suponiendo que no hay pérdidas térmicas:  
¿Cuánto calor en kcal deberá transferir la resistencia al agua para calentar todo el tanque?  

\vspace{0.5cm}

2) Sofía tiene un \emph{Freezer Horizontal Gafa 420L}. La temperatura ambiente es de $25^\circ C$, y el interior del freezer debe mantenerse a $-18^\circ C$.  
Las paredes del freezer tienen un espesor de $4cm$ y están hechas de un material con una conductividad térmica de $0.03 \frac{ J}{m \cdot s \cdot K}$.  
Si la superficie total del freezer es de $3 m^2$, ¿cuánto calor se transfiere por hora a través de las paredes del freezer?  

\vspace{0.5cm}

3)
\begin{itemize}
    \item $250^\circ K \xrightarrow{\text{convertir}}  .^\circ F $
    \item $410^\circ F \xrightarrow{\text{convertir}}  .^\circ K $
    \item $0,25^\circ BTU \xrightarrow{\text{convertir}} kcal $
    \item $850^\circ kcal \xrightarrow{\text{convertir}}  BTU $
\end{itemize}
\end{document}

   
   
